\subsection{A light at the table}

On the day his son is named, Zachary blesses God for the light that will
visit those who sit in darkness. He remembers God's oath to Abraham
and then addresses his son John, who will prepare the Lord's ways
(Luke 1:59--79). A child at supper
can understand that you need a light to find your way. On an Advent
Sunday, let someone read these few verses aloud. Ask whom Zachary is
speaking to and whose coming he expects. Even a very small child can
name someone who will need help finding the way home tonight.

Many households light a further candle each Sunday on a wreath of
green branches. The Church's \emph{Directory on Popular Piety and the
Liturgy} describes the custom especially in German-speaking lands and
North America. It sees in the growing light the long preparation by
the prophets for the coming of Christ (no.~98). If there is a wreath
on your table, light it as you hear Zachary's words. If the whole
household rarely eats together, choose the meal you do share. Someone
living alone can read the passage aloud and pray for the person whose
company they miss. Zachary brings all who sit in darkness into the
prayer he speaks over his child.

Near Christmas a crib gives children a different part to play. Let
one place Mary and Joseph and another find Luke 2:1--7. What does the
evangelist tell us that Mary did with her own hands? A child who has
helped wrap a new baby may know before the adults do. The Directory
recalls how the scene at Greccio associated with St Francis of Assisi
helped spread devotion to the crib from the thirteenth century, and
recommends reading the Gospel as the family prepares it (no.~104).
As you set out the manger, ask the child to find the detail by which
someone arriving later that night will recognize the baby (Luke 2:12).

\subsection{When the wait hurts}

Christmas can bring a loss into sharp relief. There is a name you
used to write on a parcel. This year you come to that place on the
list and have to put the pen down. When Augustine writes to Proba
about prayer, he remembers another widow. Anna had served God in
the Temple for years. When Mary and Joseph brought Jesus there,
Anna gave thanks and spoke about him to others who were waiting for
the redemption of Israel (Luke 2:36--38, Letter 130.29). Luke lets
us see both the years of her widowhood and the hour when she saw the
child. People waiting beside her heard the news from her.

You may know exactly what to ask of God. You want someone well enough
to come home. You want to see a son or daughter from whom you have
heard nothing for a long time. Speak that desire to him. Augustine recalls Paul, who asked three times for the thorn in his
flesh to leave him. The apostle could tell the Corinthians what he
had asked and what the Lord had given him to live on in his weakness
(2 Cor. 12:7--10). Augustine draws from this the freedom to ask for
relief and the courage to go on asking for God's company when the
pain stays (Letter 130.25--27). Sometimes, he writes, prayer has
more tears than words (Letter 130.20).

If it is hard even to begin, let someone else read the day's Gospel
while you listen. The Our Father has room for the person too tired to
say more than its first words. When Christ teaches us to ask for his
kingdom to come, Augustine hears us learning to want a share in what
God has promised (Letter 130.21). If you are the one who has come to
sit beside a bed, ask whether the person would like to hear a prayer
or simply have you stay. You can remain with them when the reading
is finished.
